\par\SourceLeaf{433}%
\noindent\hspace*{1.2em}Um seinen Grenzwert zu bestimmen, projizieren wir zunächst beide Elemente auf die $x,y$-Ebene. Die Flächeninhalte der Projektionen sind $Zd\sigma$, $Zd\Sigma$. Das Vorzeichen von $Z$ gibt an, ob die Grenzen die Flächen und ihre Projektionen im gleichen oder im entgegengesetzten Sinn umlaufen. Man stelle sich die Figur als Dreieck vor, dessen Projektion auf die $x,y$-Ebene die Koordinaten
\[
x,y;\qquad x+dx,\ y+dy;\qquad x+\delta x,\ y+\delta y;
\]
hat. Dann wird sein doppelter Flächeninhalt
\[
2Zd\sigma=dx.\delta y-dy.\delta x
\]
sein. Die Projektion des entsprechenden Elements auf der Kugeloberfläche hat die Koordinaten
\[
\begin{array}{cc}
X,&Y,\\[4pt]
X+\dfrac{\partial X}{\partial x}.dx+\dfrac{\partial X}{\partial y}.dy,&Y+\dfrac{\partial Y}{\partial x}.dx+\dfrac{\partial Y}{\partial y}.dy,\\[6pt]
X+\dfrac{\partial X}{\partial x}.\delta x+\dfrac{\partial X}{\partial y}.\delta y,&Y+\dfrac{\partial Y}{\partial x}.\delta x+\dfrac{\partial Y}{\partial y}.\delta y.
\end{array}
\]
Daraus ergibt sich der doppelte Flächeninhalt des Elements,
\[
\begin{aligned}
2Zd\Sigma={}&\left(\frac{\partial X}{\partial x}.dx+\frac{\partial X}{\partial y}.dy\right)\left(\frac{\partial Y}{\partial x}.\delta x+\frac{\partial Y}{\partial y}.\delta y\right)\\
&-\left(\frac{\partial X}{\partial x}.\delta x+\frac{\partial X}{\partial y}.\delta y\right)\left(\frac{\partial Y}{\partial x}.dx+\frac{\partial Y}{\partial y}.dy\right)\\
={}&\left(\frac{\partial X}{\partial x}.\frac{\partial Y}{\partial y}-\frac{\partial X}{\partial y}.\frac{\partial Y}{\partial x}\right)(dx.\delta y-dy.\delta x).
\end{aligned}
\]
Das Krümmungsmaß ist also
\[
=\frac{\partial X}{\partial x}\frac{\partial Y}{\partial y}-\frac{\partial X}{\partial y}\frac{\partial Y}{\partial x}=\omega.
\]
Nun ergibt sich, da
\[
\begin{gathered}
X=-tZ,\qquad Y=-uZ,\\
(1+tt+uu)ZZ=1:
\end{gathered}
\]
\[
\begin{aligned}
dX&=-Z^{3}(1+uu)dt+Z^{3}tu.du\\
dY&=+Z^{3}tu.dt-Z^{3}(1+tt)du,
\end{aligned}
\]
